% Double-column sn-jnl layout: use figure* for these 180-mm-wide panels.
% Place after the paragraph reporting complete-test classification results.
\begin{figure*}[t]
\centering
\includegraphics[width=\textwidth]{figures/Fig_dataset_accuracy.pdf}
\caption{\textbf{Image classification across three public datasets.}
\textbf{a}, Accuracy on the complete official test partition of each
dataset, comprising 10,000 images. Floating-point, simulated W8A8,
analytic FROST and finite-stream FROST share the same trained Transformer
weights within each model pair. W8A8 quantizes the linear and attention
products while retaining floating-point nonlinear functions.
\textbf{b}, Paired accuracy change relative to the corresponding
floating-point model, measured in percentage points; zero denotes equal
accuracy. Open markers indicate individual trained models. Filled markers
and error bars show means and sample standard deviations across three
independently trained models. Each full-test finite-stream evaluation uses
$L=256$ and one prespecified stream realization. Model and calibration
selection use separate subsets of the official training partition.
These results concern a compact Transformer trained from scratch for each
dataset.}
\label{fig:dataset_accuracy}
\end{figure*}

% Place after the paragraph interpreting the stream-length diagnostic.
\begin{figure*}[t]
\centering
\includegraphics[width=\textwidth]{figures/Fig_dataset_streams.pdf}
\caption{\textbf{Stream-length dependence on image-classification tasks.}
\textbf{a}, Accuracy change from the corresponding floating-point model
as a function of bitstream length. \textbf{b}, Relative centred-logit
error, with dotted lines showing dataset-specific mean errors for analytic
evaluation of the calibrated finite-state operators.
Each dataset uses the same prespecified, class-balanced subset of 1,000
test images at every length. The calibrated $L=256$ predictions are
extracted from the complete-test evaluation in
Fig.~\ref{fig:dataset_accuracy}. Open markers show individual trained-model
results; filled markers and error bars indicate the mean and sample
standard deviation across three trained models using stream seed 101.
Additional stream seeds at $L=256$ are reported separately in the source
data. Bitstream length excludes softmax and scalar FSM burn-in periods
of 256 and 64 cycles, respectively, and is not a measurement of network
latency.}
\label{fig:dataset_streams}
\end{figure*}

% Place after the paragraph comparing direct transfer and activation calibration.
\begin{figure*}[t]
\centering
\includegraphics[width=\textwidth]{figures/Fig_dataset_calibration.pdf}
\caption{\textbf{Activation calibration without updating Transformer weights.}
\textbf{a}, Relative centred-logit error for direct transfer of the
generic HI-SMURF operator and for activation-calibrated operators.
\textbf{b}, Disagreement with the corresponding floating-point
class predictions. All configurations evaluate the complete stochastic
network at $L=256$ using Sobol arithmetic streams and IID FSM drives.
Calibration preserves the 16-state, rank-8 operator architecture while
updating its parameters and selecting an encoding bound using separate
calibration and calibration-validation images. No task labels are used
for calibration. Comparisons use the same fixed subset of 1,000 test
images per dataset. Faint lines connect the same trained model across
conditions. Open markers denote individual models, and filled markers
and error bars denote the mean and sample standard deviation across
three trained models at stream seed 101.}
\label{fig:dataset_calibration}
\end{figure*}
